\section{Вынужденное движение}

Рассмотрим систему 2-го порядка, заданную дифференциальным уравнением:
\begin{equation}
\ddot{y} + a_1 \dot{y} + a_0 y = a_0 u
\end{equation}

С использованием блоков элементарных операций построим структурную схему данной системы. На структурной схеме отметим блоки, на которых задаются начальные условия $y(0)$, $\dot{y}(0)$. Рекомендуется использовать схему, построенную в Задании 1 Лабораторной работы 2.

Для второго, третьего и четвертого наборов коэффициентов $a_1$ и $a_0$, рассчитанных в рамках Задания 1 Лабораторной работы 2, выполним моделирование движения системы для трех случаев начальных условий:

\begin{itemize}
    \item $y(0) = -1$, $\dot{y}(0) = 0$;
    \item $y(0) = 0$, $\dot{y}(0) = 0$;
    \item $y(0) = 1$, $\dot{y}(0) = 0$.
\end{itemize}

Входные воздействия $u(t)$ возьмем из Таблицы 1 для нашего варианта 27 ($u(t)=\{2,0.7t,\cos(4t)\}$).

Для повышения наглядности рекомендуется для каждой системы и каждого входного воздействия построить графики выхода с различными начальными условиями на одних координатных осях. Таким образом, всего должно получиться по 3 изображения для каждой системы, на каждом из которых будет 3 траектории выхода, полученные для разных начальных условий.
Сделаем выводы.


Возьмём структурную схему из \textbf{Задания 1 Лабораторной работы 2} для наших входных значений  $u(t)$ (см. рис. \ref{1}-\ref{3}). 

\begin{figure}[H]
    \centering
    \includegraphics[width=1\linewidth]{/home/ns/Desktop/ITMO-courses/Linear_automatic_control_systems/lab_3/code/images/scheme_1_1.png}
    \caption{Структурная схема для $u(t)=2$}
    \label{1}
\end{figure}
\begin{figure}[H]
    \centering
    \includegraphics[width=1\linewidth]{/home/ns/Desktop/ITMO-courses/Linear_automatic_control_systems/lab_3/code/images/scheme_1_2.png}
    \caption{Структурная схема для $u(t)=0.7t$}
    \label{2}
\end{figure}
\begin{figure}[H]
    \centering
    \includegraphics[width=1\linewidth]{/home/ns/Desktop/ITMO-courses/Linear_automatic_control_systems/lab_3/code/images/scheme_1_3.png}
    \caption{Структурная схема для $u(t)=\cos(4t)$}
    \label{3}
\end{figure}


Рассмотрим второй, третий и четвертый набор коэффициентов $a_0 ,a_1$, рассчитанных в рамках \textbf{Задания 1 Лаб. работы 2}:

\begin{par}
\begin{flushleft}
\texttt{№  a\_1		 a\_0}
\end{flushleft}
\end{par}

\begin{par}
\begin{flushleft}
\texttt{2   7.400	 38.690}
\end{flushleft}
\end{par}

\begin{par}
\begin{flushleft}
\texttt{3   0.000	841.000}
\end{flushleft}
\end{par}

\begin{par}
\begin{flushleft}
\texttt{4  -1.400	 25.490}
\end{flushleft}
\end{par}

Выполним для них моделирование движения системы согласно данным нам начальным условиям и входным воздействиям и посмотрим на полученные графики:

Для $u(t)=2$ на рис. \ref{4} - \ref{42}

\begin{figure}[H]
    \centering
    \includegraphics[width=1\linewidth]{/home/ns/Desktop/ITMO-courses/Linear_automatic_control_systems/lab_3/code/images/plot_1_u2_1.png}
    \caption{Набор 2}
    \label{4}
\end{figure}
\begin{figure}[H]
    \centering
    \includegraphics[width=1\linewidth]{/home/ns/Desktop/ITMO-courses/Linear_automatic_control_systems/lab_3/code/images/plot_1_u2_2.png}
    \caption{Набор 3}
    \label{41}
\end{figure}
\begin{figure}[H]
    \centering
    \includegraphics[width=1\linewidth]{/home/ns/Desktop/ITMO-courses/Linear_automatic_control_systems/lab_3/code/images/plot_1_u2_3.png}
    \caption{Набор 4}
    \label{42}
\end{figure}
Для $u(t)=0.7t$ на рис. \ref{5} - \ref{52}

\begin{figure}[H]
    \centering
    \includegraphics[width=1\linewidth]{/home/ns/Desktop/ITMO-courses/Linear_automatic_control_systems/lab_3/code/images/plot_1_u07t_1.png}
    \caption{Набор 2}
    \label{5}
\end{figure}
\begin{figure}[H]
    \centering
    \includegraphics[width=1\linewidth]{/home/ns/Desktop/ITMO-courses/Linear_automatic_control_systems/lab_3/code/images/plot_1_u07t_2.png}
    \caption{Набор 3}
    \label{51}
\end{figure}
\begin{figure}[H]
    \centering
    \includegraphics[width=1\linewidth]{/home/ns/Desktop/ITMO-courses/Linear_automatic_control_systems/lab_3/code/images/plot_1_u07t_3.png}
    \caption{Набор 4}
    \label{52}
\end{figure}
Для $u(t)=\cos(4t)$ на рис. \ref{6} - \ref{62}

\begin{figure}[H]
    \centering
    \includegraphics[width=1\linewidth]{/home/ns/Desktop/ITMO-courses/Linear_automatic_control_systems/lab_3/code/images/plot_1_ucos4t_1.png}
    \caption{Набор 2}
    \label{6}
\end{figure}
\begin{figure}[H]
    \centering
    \includegraphics[width=1\linewidth]{/home/ns/Desktop/ITMO-courses/Linear_automatic_control_systems/lab_3/code/images/plot_1_ucos4t_2.png}
    \caption{Набор 3}
    \label{61}
\end{figure}
\begin{figure}[H]
    \centering
    \includegraphics[width=1\linewidth]{/home/ns/Desktop/ITMO-courses/Linear_automatic_control_systems/lab_3/code/images/plot_1_ucos4t_3.png}
    \caption{Набор 4}
    \label{62}
\end{figure}


\textbf{Анализ и выводы}

При каждом из входных воздействий мы видим, что для любых взятых нами начальных условий 2-й набор коэффициентов делает систему ассимптотически устойчивой, 3-й - граница устойчивости и  4-й - делает систему неуйстойчивой. На рисунках видно, что при ассимптотической устойчивости графики систем со временем сходятся к функции входного воздействия, при границе усточивости колеблются в некоторой окрестности и при неустойчивости убегают от вида функции входного воздействия. Таким образом, в данном задании мы довольно наглядно на эксперименте проверили выводы из предыдущей лабораторной работы о том, при каких коэффициентах какой устойчивостью система обладает.


\section{Качество переходных процессов}


Для системы 3-го порядка, заданной передаточной функцией
$$W(s)=\frac{|\lambda_1 \lambda_2 \lambda_3 |}{(s-\lambda_1 )(s-\lambda_2 )(s-\lambda_3 )},$$
исследуем зависимости качества переходной характеристики от выбора полюсов передаточной функции, для оценки качества используем такие показатели, как \textbf{перерегулирование} и \textbf{время переходного процесса}. Для этого зададимся \textbf{шестью} наборами полюсов $\lambda_1 ,\lambda_2$ и $\lambda_3$ с отрицательной вещественной частью. Половину наборов возьмем чисто вещественной, а во вторую половину включим комплексно-сопряженные полюса. В ходе исследования рекомендуется начинать с корней, имеющих равные вещественные части, а затем, фиксируя один или пару корней из набора, варьировать остальные. Проведем моделирование для выбранных наборов полюсов, определим значения времени переходных процессов и перерегулирования для каждого случая, сопоставим полученные результаты и сделаем выводы. Также сравним полученные значения перерегулирования с их верхними оценками на основании степени колебательности.

Выберем следующие наборы полюсов $\lambda_1 ,\lambda_2 ,\lambda_2$:
\\
$ \lbrace -1,-1,-1\rbrace ,\\\lbrace -1,-1,-3\rbrace ,\\\lbrace -1,-3,-5\rbrace ,\\\lbrace -1,-1+4i,-1-4i\rbrace ,\\\lbrace -4,-1+4i,-1-4i\rbrace ,\\\lbrace -4,-1+8i,-1-8i\rbrace$


Их расположение на комплексной плоскости на рисунках \ref{7}-\ref{12}.

\begin{figure}[H]
    \centering
    \includegraphics[width=0.6\linewidth]{/home/ns/Desktop/ITMO-courses/Linear_automatic_control_systems/lab_3/code/images/poles_1.png}
    \caption{Набор полюсов 1}
    \label{7}
\end{figure}


\begin{figure}[H]
    \centering
    \includegraphics[width=0.6\linewidth]{/home/ns/Desktop/ITMO-courses/Linear_automatic_control_systems/lab_3/code/images/poles_2.png}
    \caption{Набор полюсов 2}
    \label{8}
\end{figure}


\begin{figure}[H]
    \centering
    \includegraphics[width=0.6\linewidth]{/home/ns/Desktop/ITMO-courses/Linear_automatic_control_systems/lab_3/code/images/poles_3.png}
    \caption{Набор полюсов 3}
    \label{9}
\end{figure}


\begin{figure}[H]
    \centering
    \includegraphics[width=0.6\linewidth]{/home/ns/Desktop/ITMO-courses/Linear_automatic_control_systems/lab_3/code/images/poles_4.png}
    \caption{Набор полюсов 4}
    \label{10}
\end{figure}


\begin{figure}[H]
    \centering
    \includegraphics[width=0.6\linewidth]{/home/ns/Desktop/ITMO-courses/Linear_automatic_control_systems/lab_3/code/images/poles_5.png}
    \caption{Набор полюсов 5}
    \label{11}
\end{figure}


\begin{figure}[H]
    \centering
    \includegraphics[width=0.6\linewidth]{/home/ns/Desktop/ITMO-courses/Linear_automatic_control_systems/lab_3/code/images/poles_6.png}
    \caption{Набор полюсов 6}
    \label{12}
\end{figure}


Построим графики переходных процессов (см. рис. \ref{13}-\ref{18}).

\begin{figure}[H]
    \centering
    \includegraphics[width=1\linewidth]{/home/ns/Desktop/ITMO-courses/Linear_automatic_control_systems/lab_3/code/images/per_proc_1.png}
    \caption{График переходного процесса при наборе полюсов 1}
    \label{13}
\end{figure}


\begin{figure}[H]
    \centering
    \includegraphics[width=1\linewidth]{/home/ns/Desktop/ITMO-courses/Linear_automatic_control_systems/lab_3/code/images/per_proc_2.png}
    \caption{График переходного процесса при наборе полюсов 2}
    \label{14}
\end{figure}


\begin{figure}[H]
    \centering
    \includegraphics[width=1\linewidth]{/home/ns/Desktop/ITMO-courses/Linear_automatic_control_systems/lab_3/code/images/per_proc_3.png}
    \caption{График переходного процесса при наборе полюсов 3}
    \label{15}
\end{figure}


\begin{figure}[H]
    \centering
    \includegraphics[width=1\linewidth]{/home/ns/Desktop/ITMO-courses/Linear_automatic_control_systems/lab_3/code/images/per_proc_4.png}
    \caption{График переходного процесса при наборе полюсов 4}
    \label{16}
\end{figure}


\begin{figure}[H]
    \centering
    \includegraphics[width=1\linewidth]{/home/ns/Desktop/ITMO-courses/Linear_automatic_control_systems/lab_3/code/images/per_proc_5.png}
    \caption{График переходного процесса при наборе полюсов 5}
    \label{17}
\end{figure}


\begin{figure}[H]
    \centering
    \includegraphics[width=1\linewidth]{/home/ns/Desktop/ITMO-courses/Linear_automatic_control_systems/lab_3/code/images/per_proc_6.png}
    \caption{График переходного процесса при наборе полюсов 6}
    \label{18}
\end{figure}


\textbf{Анализ и сравнение графиков}


Для начала введём основным понятия, необходимые для анализа.

\textit{Установившееся значение и перерегулирование} относятся к динамическим показателям качества, то есть, они оцениваются по переходной характеристике.


Установившееся значение $y_{\textrm{уст}}$ - это предел функции $y(t)$ при $t\to \infty$. Перерегулирование - это временный выброс сигнала над $y_{\textrm{уст}}$, рассчитывается по формуле: $\sigma =\frac{|y_{max} -y_{уст} |}{|y_{уст} |}\cdot 100\%$. Время переходного процесса $\Delta$ - это интервал времени от начала внешнего воздействия до момента, после которого разница между сигналом и $y_{\textrm{уст}}$ не превышает определенного значения, в данной л/р примем $\Delta =0.05y_{\textrm{уст}}$: $|y(t)-y_{\textrm{уст}} |<\Delta$.

\textit{Степень колебательности} относится к корневым показателям качества, они определятся по расположению полюсов на комплексной плоскости. Степень колебательности равна $\mu =max_k \left|\frac{\textrm{Im}(\lambda_k )}{\textrm{Re}(\lambda_k )}\right|,~~k=\overline{1,n}$ для $n$ полюсов системы. Доминирующими полюсами называются те, что расположены ближе к мнимой оси. В случае, когда пара комплексных полюсов доминирует, для оценки перерегулирования можно использовать формулу: $\sigma_{\%} \le e^{-\pi /\mu } \cdot 100\%$.

Для наших наборов получим:

\begin{itemize}
\setlength{\itemsep}{-1ex}
   \item{\begin{flushleft} $\lbrace -1,-1,-1\rbrace$: перерегулирование $0\%$, время переходного процесса $\approx 6.296$ \end{flushleft}}
   \item{\begin{flushleft} $\lbrace -1,-1,-3\rbrace$: перерегулирование $0\%$, время переходного процесса $\approx 5.13$ \end{flushleft}}
   \item{\begin{flushleft} $\lbrace -1,-3,-5\rbrace$: перерегулирование $0\%$, время переходного процесса $\approx 3.624$ \end{flushleft}}
   \item{\begin{flushleft} $\lbrace -1,-1+4i,-1-4i\rbrace V$: перерегулирование $0\%$, время переходного процесса $\approx 2.781$ \end{flushleft}}
   \item{\begin{flushleft} $\lbrace -4,-1+4i,-1-4i\rbrace$: перерегулирование $27.8\%$, время переходного процесса $\approx 2.739$ \end{flushleft}}
   \item{\begin{flushleft} $\lbrace -4,-1+8i,-1-8i\rbrace$: перерегулирование $17.5\%$, время переходного процесса $\approx 2.176$ \end{flushleft}}
\end{itemize}


Для всех наборов с чисто вещественными корнями степень колебательности равна нулю, для остальных:


\begin{itemize}
\setlength{\itemsep}{-1ex}
   \item{\begin{flushleft} $\lbrace -1,-1+4i,-1-4i\rbrace$: $\mu =4$ \end{flushleft}}
   \item{\begin{flushleft} $\lbrace -4,-1+4i,-1-4i\rbrace$: $\mu =4$ \end{flushleft}}
   \item{\begin{flushleft} $\lbrace -4,-1+8i,-1-8i\rbrace$: $\mu =8$ \end{flushleft}}
\end{itemize}


Для последних двух наборов, получивших перерегулирование сделаем его верхнюю оценку на основе колебательности и сравним с полученными результатами по явной формуле для перерегулирования:


\begin{itemize}
\setlength{\itemsep}{-1ex}
   \item{\begin{flushleft} $\lbrace -4,-1+4i,-1-4i\rbrace$: $\sigma_{\%} \le e^{-\pi /4} \cdot 100\%\approx 45.6$ \end{flushleft}}
   \item{\begin{flushleft} $\lbrace -4,-1+8i,-1-8i\rbrace$: $\sigma_{\%} \le e^{-\pi /8} \cdot 100\%\approx 67.5$ \end{flushleft}}
\end{itemize}

Оценка получилась довольно грубой, но она выполняется.


\textbf{Анализ и выводы}


\begin{par}
\begin{flushleft}
\&\&\&ДОПИСАТЬ\&\&\&
\end{flushleft}
\end{par}




\section{Приложение}

\matlabtitle{Задание 1}
\begin{matlabcode}
a1_a0_arr = [ [7.4 38.69]; [0 841]; [-1.4 25.49] ];
y0_dot_y0_arr = [ [-1 0]; [0 0]; [1 0] ];

simtime = 4;

for i=1:3
    a_1 = a1_a0_arr(i,1);
    a_0 = a1_a0_arr(i,2);
    for k=1:3
        y_0 = y0_dot_y0_arr(k,1);
        dot_y_0 = y0_dot_y0_arr(k,2);
    
        plot_1 = sim("lab3_1_1.slx", simtime);
        plot(plot_1.y.Time, plot_1.y.Data, LineWidth=1, LineStyle="-", ...
            DisplayName=sprintf('$y(0)=%d, \\dot{y}(0)=%d$', y_0, dot_y_0));
        hold on;
    
    end
    grid on;
    legend(Location="best", FontSize=14, Interpreter="latex");
    xlabel("$t, c$", Interpreter="latex");
    ylabel("$y(t)$", Interpreter="latex");
    set(gcf, Position=[100, 100, 1200, 600]);   
    set(gca, FontSize=14);
    
    hold off;
    filename = sprintf('images/plot_1_u2_%d.png', i);
    exportgraphics(gcf, filename, 'Resolution', 200)
end

for i=1:3
    a_1 = a1_a0_arr(i,1);
    a_0 = a1_a0_arr(i,2);
    for k=1:3
        y_0 = y0_dot_y0_arr(k,1);
        dot_y_0 = y0_dot_y0_arr(k,2);
    
        plot_1 = sim("lab3_1_2.slx", simtime);
        plot(plot_1.y.Time, plot_1.y.Data, LineWidth=1, LineStyle="-", ...
            DisplayName=sprintf('$y(0)=%d, \\dot{y}(0)=%d$', y_0, dot_y_0));
        hold on;
    
    end
    grid on;
    legend(Location="best", FontSize=14, Interpreter="latex");
    xlabel("$t, c$", Interpreter="latex");
    ylabel("$y(t)$", Interpreter="latex");
    set(gcf, Position=[100, 100, 1200, 600]);   
    set(gca, FontSize=14);
    
    hold off;
    filename = sprintf('images/plot_1_u07t_%d.png', i);
    exportgraphics(gcf, filename, 'Resolution', 200)
end

for i=1:3
    a_1 = a1_a0_arr(i,1);
    a_0 = a1_a0_arr(i,2);
    for k=1:3
        y_0 = y0_dot_y0_arr(k,1);
        dot_y_0 = y0_dot_y0_arr(k,2);
    
        plot_1 = sim("lab3_1_3.slx", simtime);
        plot(plot_1.y.Time, plot_1.y.Data, LineWidth=1, LineStyle="-", ...
            DisplayName=sprintf('$y(0)=%d, \\dot{y}(0)=%d$', y_0, dot_y_0));
        hold on;
    
    end
    grid on;
    legend(Location="best", FontSize=14, Interpreter="latex");
    xlabel("$t, c$", Interpreter="latex");
    ylabel("$y(t)$", Interpreter="latex");
    set(gcf, Position=[100, 100, 1200, 600]);   
    set(gca, FontSize=14);
    
    hold off;
    filename = sprintf('images/plot_1_ucos4t_%d.png', i);
    exportgraphics(gcf, filename, 'Resolution', 200)
end
\end{matlabcode}

\matlabtitle{Задание 2}
\begin{matlabcode}
pole_sets = [
    [-1, -1, -1];
    [-1, -1, -3];
    [-1, -3, -5];
    [-1, -1+4i, -1-4i];
    [-4, -1+4i, -1-4i];
    [-4, -1+8i, -1-8i]
];


for i = 1:6
    current_poles = pole_sets(i,:);
    
    figure;
    xline(0, color="black", linestyle='-', LineWidth=0.5, ...
        HandleVisibility='off');
    hold on;
    yline(0, color="black", linestyle='-', LineWidth=0.5, ...
        HandleVisibility='off');
    hold on;
    pole_strs = arrayfun(@(p) sprintf('%.1f%+.1fi', real(p), imag(p)), ...
        current_poles, UniformOutput=false);
    p1 = plot(real(current_poles(1)), imag(current_poles(1)), 'g.', ...
        MarkerSize=20, ...
        DisplayName=sprintf('$\\lambda_{1}=%s$', pole_strs{1}));
    hold on;
    p2 = plot(real(current_poles(2)), imag(current_poles(2)), 'b.', ...
        MarkerSize=20, ...
        DisplayName=sprintf('$\\lambda_{2}=%s$', pole_strs{2}));
    hold on;
    p3 = plot(real(current_poles(3)), imag(current_poles(3)), 'r.', ...
        MarkerSize=20, ...
        DisplayName=sprintf('$\\lambda_{3}=%s$', pole_strs{3}));
  
    grid on;
    xlabel("Re\{z\}", Interpreter="latex");
    ylabel("Im\{z\}", Interpreter="latex");
    title(sprintf('Set %d', i), Interpreter="latex");
    legend(Location="best", FontSize=12, Interpreter="latex");
    set(gca, FontSize=12);

    xlim([-5, 1.0]);
    ylim([-8.5, 8.5]);
    hold off;
    
    filename = sprintf('images/poles_%d.png', i);
    exportgraphics(gcf, filename, 'Resolution', 200);
end

t = 0:0.001:10; 
u = ones(size(t));

for i = 1:6
    lambda_1 = pole_sets(i,1);
    lambda_2 = pole_sets(i,2);
    lambda_3 = pole_sets(i,3);

    W = zpk([], [lambda_1, lambda_2, lambda_3], ...
        abs(lambda_1*lambda_2*lambda_3));
    y = lsim(W, u, t);
    y_st = dcgain(W);

    delta = 0.05 * y_st;
    transition_time = NaN;
    
    for j = 1:length(t)
        if abs(y(j) - y_st) < delta
            if all(abs(y(j:end) - y_st) < delta)
                transition_time = t(j);
                break;
            end
        end
    end
    
    fprintf('Set %d: Время переходного процесса = %.3f с\n', ...
        i, transition_time);
    
    figure;
    yline(y_st, 'r--', 'LineWidth', 1, DisplayName="$y_{steady}$");
    hold on;
    plot(t, y, 'LineWidth', 2, 'LineStyle', '-', DisplayName="$y(t)$");
    grid on;
    ylabel("$y(t)$", Interpreter="latex");
    xlabel("$t, c$", Interpreter="latex");
    title(sprintf('Set %d', i), Interpreter="latex");
    legend(Location="best", FontSize=12, Interpreter="latex");
    set(gcf, Position=[100, 100, 1200, 600]);
    set(gca, FontSize=12);
    ylim([0, 1.4]);
    filename = sprintf('images/per_proc_%d.png', i);
    exportgraphics(gcf, filename, 'Resolution', 200);
end

for i = 5:6
    lambda_1 = pole_sets(i,1);
    lambda_2 = pole_sets(i,2);
    lambda_3 = pole_sets(i,3);

    W = zpk([], [lambda_1, lambda_2, lambda_3], ...
        abs(lambda_1*lambda_2*lambda_3));
    y = lsim(W, u, t);
    y_st = dcgain(W);
    fprintf('Set %d: Перерегулирование = %.4f\n %', ...
        i, abs(max(y)-y_st)/y_st * 100);
end


\end{matlabcode}
